
\Problem{Vectors de bits harmoniosos}

\Statement

Diem que un vector $v$ de $n$ bits és $d$-harmoniós per a un natural $d$ si,
per a cada $i$ amb $0\le i<n$, la diferència
(en valor absolut) entre el nombre
de zeros i el nombre de uns en $v[0,\dots,i]$ és inferior
o igual a $d$.

\medskip

Per exemple,
hi ha 4 vectors
1-harmoniosos de llargada 4:
0101,
0110,
1001, i
1010.
Igualment,
hi ha 12 vectors
2-harmoniosos de llargada 4:
0010,
0011,
0100,
0101,
0110,
0111,
1000,
1001,
1010,
1011,
1100 i
1101.
També, hi ha 4 vectors 1-harmoniosos de llargada 3:
010,
011,
100 i
101. Fixeu-vos que no hi ha vectors
0-harmoniosos per a $n>0$, però que n'hi ha un de llargada $n=0$ (el vector buit).

\medskip

Escriviu un programa que, per a uns $n$ i $d$ donats escrigui quants vectors
de mida $n$ són $d$-harmoniosos.

\Input

L'entrada consisteix en una seqüència de parells naturals $n, d \ge 0$.


\Output

Per a cada $n$ i cada $d$, escriviu quants vectors
de mida $n$ són $d$-harmoniosos.


\Sample
